% Figure: dynamic amplification factor of the forced SDOF oscillator
% versus frequency ratio, for several damping ratios.
% X/(F0/k) = 1/sqrt((1-r^2)^2 + (2 zeta r)^2), r = omega/omega_n.
\begin{figure}[htbp]
\centering
\begin{tikzpicture}
\begin{axis}[
  width=12cm, height=7.5cm,
  xlabel={$r=\omega/\omega_n$}, ylabel={$X\,/\,(F_0/k)$},
  xmin=0, xmax=2.5, ymin=0, ymax=10.5,
  axis lines=left,
  domain=0.001:2.5, samples=400,
  legend pos=north east,
  legend cell align=left,
  every axis x label/.style={at={(current axis.right of origin)}, anchor=west},
  every axis y label/.style={at={(current axis.above origin)}, anchor=south},
  grid=major, grid style={enggray!20},
  tick label style={font=\scriptsize},
  label style={font=\small},
  legend style={font=\scriptsize},
]
% zeta = 0.05 -> peak ~ 1/(2 zeta) = 10
\addplot[engred, very thick] {1/sqrt((1-x^2)^2 + (2*0.05*x)^2)};
\addlegendentry{$\zeta=0.05$ ($Q=10$)}
% zeta = 0.10
\addplot[engorange, very thick] {1/sqrt((1-x^2)^2 + (2*0.10*x)^2)};
\addlegendentry{$\zeta=0.10$ ($Q=5$)}
% zeta = 0.25
\addplot[engblue, very thick] {1/sqrt((1-x^2)^2 + (2*0.25*x)^2)};
\addlegendentry{$\zeta=0.25$}
% zeta = 0.50
\addplot[engamber, very thick, dashed] {1/sqrt((1-x^2)^2 + (2*0.50*x)^2)};
\addlegendentry{$\zeta=0.50$}
% static deflection line
\addplot[enggray, thin, forget plot] {1};
\node[enggray, font=\scriptsize, anchor=south west] at (axis cs:0.04,1.05) {static deflection};
\end{axis}
\end{tikzpicture}
\caption[Dynamic amplification of the forced oscillator]{Steady-state
amplitude of $m\ddot{x}+c\dot{x}+kx=F_0\cos(\omega t)$, normalised to
the static deflection $F_0/k$, against the frequency ratio
$r=\omega/\omega_n$, for four damping ratios. Light damping
($\zeta=0.05$, red) produces a sharp peak of height
$\approx 1/(2\zeta)=10$ near $r=1$. Increasing damping lowers and
broadens the peak. Below $r\approx 0.3$ every curve sits near the
static deflection: the system follows the force quasi-statically.
Above $r\approx 2$ the response rolls off as $r^{-2}$ regardless of
damping: mass inertia dominates. The peak occurs not exactly at $r=1$
but at $r=\sqrt{1-2\zeta^2}$, and disappears once $\zeta>1/\sqrt2$.
Computed from $X/(F_0/k)=[(1-r^2)^2+(2\zeta r)^2]^{-1/2}$.}
\label{fig:vol03ch06:frequency-response}
\end{figure}
